% --- Archon physics formalization source begin ---
% archon:physics
% archon:covers PhyXMiniProblems/problem_phyx_mini_0925.lean
% archon:source-report reports/phyx_mini/problem_phyx_mini_0925.source.json

\chapter{Physics problem phyx\_mini\_0925}
\label{ch:PhyXMiniProblems_problem_phyx_mini_0925}

\paragraph{Problem source.}
\#\# Physical scenario

Figure shows a 2.0-cm-diameter solenoid passing through the center of a 6.0-cm-diameter loop. The magnetic field inside the solenoid is 0.20 T.

\#\# Question

What is the magnetic flux through the loop when it is tilted at a 60\textbackslash{}(\textasciicircum{}\{\textbackslash{}circ\}\textbackslash{}) angle?

\#\# Answer choices

- A: A: \textbackslash{}( 6.3 \textbackslash{}times 10\textasciicircum{}\{-6\} \textbackslash{}text\{ Wb\} \textbackslash{})

- B: B: \textbackslash{}( 6.3 \textbackslash{}times 10\textasciicircum{}\{-5\} \textbackslash{}text\{ Wb\} \textbackslash{})

- C: C: \textbackslash{}( 6.3 \textbackslash{}times 10\textasciicircum{}\{-4\} \textbackslash{}text\{ Wb\} \textbackslash{})

- D: D: \textbackslash{}( 6.3 \textbackslash{}times 10\textasciicircum{}\{-7\} \textbackslash{}text\{ Wb\} \textbackslash{})

\#\# Figure caption (auxiliary; use the image as primary evidence)

The image shows a mechanical setup with the following elements:

1. A horizontal rectangular block with vertical lines, possibly indicating texture or material properties.
2. Two cylindrical rods placed underneath the block at an angle.
3. The rods are orange and positioned to support the block.
4. There is an angle marker indicating 60 degrees between the block and one of the rods.
5. The angle is marked with an arc and labeled "60°," indicating the inclination of one rod relative to the perpendicular of the block.

\#\# Dataset metadata

Electromagnetic Induction; Physical Model Grounding Reasoning, Spatial Relation Reasoning

\paragraph{Recorded answer/context.}
B: B: \textbackslash{}( 6.3 \textbackslash{}times 10\textasciicircum{}\{-5\} \textbackslash{}text\{ Wb\} \textbackslash{})

\paragraph{Figure/image path.}
/root/proposal\_for\_physic/science-mango/phyx\_mini\_run/phyx\_data/test\_image/925.png

\paragraph{Formalization target.}
create a compiling Lean file with sorry bodies at `PhyXMiniProblems/problem\_phyx\_mini\_0925.lean`. The Lean declarations must preserve the physical quantities, dimensions or dimensional roles, figure labels, governing-law hypotheses, and final relation expressed by this problem.
Use Mathlib/Physlib names found through LeanExplore where available. If a physics API is missing, introduce faithful local abstractions rather than scalar placeholder aliases.

\begin{theorem}[Physics formalization target]
\label{thm:physics:phyx_mini_0925:target}
\lean{PhyXMiniProblems.ProblemPhyXMini0925.problem_phyx_mini_0925}
\uses{def:physics:phyx-mini-0925:phyxminiproblems-problemphyxmini0925-magneticfluxinwebers, def:physics:phyx-mini-0925:phyxminiproblems-problemphyxmini0925-tiltedsolenoidloopsetup, def:physics:phyx-mini-0925:phyxminiproblems-problemphyxmini0925-matchesproblemdescription, def:physics:phyx-mini-0925:phyxminiproblems-problemphyxmini0925-matchesprimarytiltfigure, def:physics:phyx-mini-0925:phyxminiproblems-problemphyxmini0925-hasphysicalsolenoidloopparameters, def:physics:phyx-mini-0925:phyxminiproblems-problemphyxmini0925-satisfiescircularareageometry, def:physics:phyx-mini-0925:phyxminiproblems-problemphyxmini0925-satisfiestiltedthreadinggeometry, def:physics:phyx-mini-0925:phyxminiproblems-problemphyxmini0925-modelsuniformconfinedsolenoidfield, def:physics:phyx-mini-0925:phyxminiproblems-problemphyxmini0925-satisfiesuniformmagneticsurfacefluxlaw, def:physics:phyx-mini-0925:phyxminiproblems-problemphyxmini0925-roundstodisplayedflux, def:physics:phyx-mini-0925:phyxminiproblems-problemphyxmini0925-isclosestdisplayedchoice}
The assigned autoformalize agent should translate this physics problem into Lean declarations in the covered file, with theorem and lemma proof bodies written as `by sorry`.
\end{theorem}
\begin{proof}
This is an autoformalization task, not a proof task. Produce statements that can later be proved by the physics prover without weakening the physical model.
\end{proof}

% --- Archon declaration topology begin ---
\section{Declaration topology}
\begin{definition}[magnetic Flux Density Dimension]
  \label{def:physics:phyx-mini-0925:phyxminiproblems-problemphyxmini0925-magneticfluxdensitydimension}
  \lean{PhyXMiniProblems.ProblemPhyXMini0925.magneticFluxDensityDimension}
  Magnetic flux density (tesla) has physical dimension M T⁻¹ C⁻¹
\end{definition}
\begin{proof}
  This is the declared model definition of magnetic Flux Density Dimension.
\end{proof}
\begin{definition}[magnetic Flux Dimension]
  \label{def:physics:phyx-mini-0925:phyxminiproblems-problemphyxmini0925-magneticfluxdimension}
  \lean{PhyXMiniProblems.ProblemPhyXMini0925.magneticFluxDimension}
  \uses{def:physics:phyx-mini-0925:phyxminiproblems-problemphyxmini0925-magneticfluxdensitydimension}
  Magnetic flux (weber) has physical dimension M L² T⁻¹ C⁻¹
\end{definition}
\begin{proof}
  This is the declared model definition of magnetic Flux Dimension.
\end{proof}
\begin{definition}[Length Magnitude]
  \label{def:physics:phyx-mini-0925:phyxminiproblems-problemphyxmini0925-lengthmagnitude}
  \lean{PhyXMiniProblems.ProblemPhyXMini0925.LengthMagnitude}
  A nonnegative, unit-independent physical length
\end{definition}
\begin{proof}
  This is the declared model definition of Length Magnitude.
\end{proof}
\begin{definition}[Magnetic Flux Density Magnitude]
  \label{def:physics:phyx-mini-0925:phyxminiproblems-problemphyxmini0925-magneticfluxdensitymagnitude}
  \lean{PhyXMiniProblems.ProblemPhyXMini0925.MagneticFluxDensityMagnitude}
  \uses{def:physics:phyx-mini-0925:phyxminiproblems-problemphyxmini0925-magneticfluxdensitydimension}
  A nonnegative, unit-independent magnetic-flux-density magnitude
\end{definition}
\begin{proof}
  This is the declared model definition of Magnetic Flux Density Magnitude.
\end{proof}
\begin{definition}[Signed Magnetic Flux]
  \label{def:physics:phyx-mini-0925:phyxminiproblems-problemphyxmini0925-signedmagneticflux}
  \lean{PhyXMiniProblems.ProblemPhyXMini0925.SignedMagneticFlux}
  \uses{def:physics:phyx-mini-0925:phyxminiproblems-problemphyxmini0925-magneticfluxdimension}
  Signed magnetic flux relative to the selected loop area normal
\end{definition}
\begin{proof}
  This is the declared model definition of Signed Magnetic Flux.
\end{proof}
\begin{definition}[nonnegative SIReadout]
  \label{def:physics:phyx-mini-0925:phyxminiproblems-problemphyxmini0925-nonnegativesireadout}
  \lean{PhyXMiniProblems.ProblemPhyXMini0925.nonnegativeSIReadout}
  Coherent-SI readout of a nonnegative dimensionful quantity
\end{definition}
\begin{proof}
  This is the declared model definition of nonnegative SIReadout.
\end{proof}
\begin{definition}[length In Meters]
  \label{def:physics:phyx-mini-0925:phyxminiproblems-problemphyxmini0925-lengthinmeters}
  \lean{PhyXMiniProblems.ProblemPhyXMini0925.lengthInMeters}
  \uses{def:physics:phyx-mini-0925:phyxminiproblems-problemphyxmini0925-lengthmagnitude, def:physics:phyx-mini-0925:phyxminiproblems-problemphyxmini0925-nonnegativesireadout}
  Read a physical length in coherent-SI metres
\end{definition}
\begin{proof}
  This is the declared model definition of length In Meters.
\end{proof}
\begin{definition}[length In Centimeters]
  \label{def:physics:phyx-mini-0925:phyxminiproblems-problemphyxmini0925-lengthincentimeters}
  \lean{PhyXMiniProblems.ProblemPhyXMini0925.lengthInCentimeters}
  \uses{def:physics:phyx-mini-0925:phyxminiproblems-problemphyxmini0925-lengthmagnitude, def:physics:phyx-mini-0925:phyxminiproblems-problemphyxmini0925-lengthinmeters}
  Read a physical length in centimetres
\end{definition}
\begin{proof}
  This is the declared model definition of length In Centimeters.
\end{proof}
\begin{definition}[area In Square Meters]
  \label{def:physics:phyx-mini-0925:phyxminiproblems-problemphyxmini0925-areainsquaremeters}
  \lean{PhyXMiniProblems.ProblemPhyXMini0925.areaInSquareMeters}
  \uses{def:physics:phyx-mini-0925:phyxminiproblems-problemphyxmini0925-nonnegativesireadout}
  Read a physical area in coherent-SI square metres
\end{definition}
\begin{proof}
  This is the declared model definition of area In Square Meters.
\end{proof}
\begin{definition}[magnetic Flux Density In Teslas]
  \label{def:physics:phyx-mini-0925:phyxminiproblems-problemphyxmini0925-magneticfluxdensityinteslas}
  \lean{PhyXMiniProblems.ProblemPhyXMini0925.magneticFluxDensityInTeslas}
  \uses{def:physics:phyx-mini-0925:phyxminiproblems-problemphyxmini0925-magneticfluxdensitymagnitude, def:physics:phyx-mini-0925:phyxminiproblems-problemphyxmini0925-nonnegativesireadout}
  Read magnetic flux density in coherent-SI teslas
\end{definition}
\begin{proof}
  This is the declared model definition of magnetic Flux Density In Teslas.
\end{proof}
\begin{definition}[magnetic Flux In Webers]
  \label{def:physics:phyx-mini-0925:phyxminiproblems-problemphyxmini0925-magneticfluxinwebers}
  \lean{PhyXMiniProblems.ProblemPhyXMini0925.magneticFluxInWebers}
  \uses{def:physics:phyx-mini-0925:phyxminiproblems-problemphyxmini0925-signedmagneticflux}
  Read signed magnetic flux in coherent-SI webers
\end{definition}
\begin{proof}
  This is the declared model definition of magnetic Flux In Webers.
\end{proof}
\begin{definition}[degrees To Radians]
  \label{def:physics:phyx-mini-0925:phyxminiproblems-problemphyxmini0925-degreestoradians}
  \lean{PhyXMiniProblems.ProblemPhyXMini0925.degreesToRadians}
  Convert the dimensionless degree readout used in the figure to radians
\end{definition}
\begin{proof}
  This is the declared model definition of degrees To Radians.
\end{proof}
\begin{definition}[Visible Loop Section]
  \label{def:physics:phyx-mini-0925:phyxminiproblems-problemphyxmini0925-visibleloopsection}
  \lean{PhyXMiniProblems.ProblemPhyXMini0925.VisibleLoopSection}
  The two pieces of the orange loop visible on opposite sides of the body
\end{definition}
\begin{proof}
  This is the declared model definition of Visible Loop Section.
\end{proof}
\begin{definition}[Tilt Angle Reference]
  \label{def:physics:phyx-mini-0925:phyxminiproblems-problemphyxmini0925-tiltanglereference}
  \lean{PhyXMiniProblems.ProblemPhyXMini0925.TiltAngleReference}
  The line from which the raster's marked tilt angle is measured
\end{definition}
\begin{proof}
  This is the declared model definition of Tilt Angle Reference.
\end{proof}
\begin{definition}[Loop Topology]
  \label{def:physics:phyx-mini-0925:phyxminiproblems-problemphyxmini0925-looptopology}
  \lean{PhyXMiniProblems.ProblemPhyXMini0925.LoopTopology}
  Topology asserted for the larger conductor
\end{definition}
\begin{proof}
  This is the declared model definition of Loop Topology.
\end{proof}
\begin{definition}[Solenoid Field Support]
  \label{def:physics:phyx-mini-0925:phyxminiproblems-problemphyxmini0925-solenoidfieldsupport}
  \lean{PhyXMiniProblems.ProblemPhyXMini0925.SolenoidFieldSupport}
  Idealized support of the solenoid field used by this problem
\end{definition}
\begin{proof}
  This is the declared model definition of Solenoid Field Support.
\end{proof}
\begin{definition}[Tilted Solenoid Loop Figure]
  \label{def:physics:phyx-mini-0925:phyxminiproblems-problemphyxmini0925-tiltedsolenoidloopfigure}
  \lean{PhyXMiniProblems.ProblemPhyXMini0925.TiltedSolenoidLoopFigure}
  \uses{def:physics:phyx-mini-0925:phyxminiproblems-problemphyxmini0925-visibleloopsection, def:physics:phyx-mini-0925:phyxminiproblems-problemphyxmini0925-tiltanglereference}
  Literal, typed content read from image 925.png. The numerical angle is a dimensionless displayed readout rather than a physical magnetic quantity
\end{definition}
\begin{proof}
  This is the declared model definition of Tilted Solenoid Loop Figure.
\end{proof}
\begin{definition}[Tilted Solenoid Loop Setup]
  \label{def:physics:phyx-mini-0925:phyxminiproblems-problemphyxmini0925-tiltedsolenoidloopsetup}
  \lean{PhyXMiniProblems.ProblemPhyXMini0925.TiltedSolenoidLoopSetup}
  \uses{def:physics:phyx-mini-0925:phyxminiproblems-problemphyxmini0925-lengthmagnitude, def:physics:phyx-mini-0925:phyxminiproblems-problemphyxmini0925-magneticfluxdensitymagnitude, def:physics:phyx-mini-0925:phyxminiproblems-problemphyxmini0925-signedmagneticflux, def:physics:phyx-mini-0925:phyxminiproblems-problemphyxmini0925-looptopology, def:physics:phyx-mini-0925:phyxminiproblems-problemphyxmini0925-solenoidfieldsupport, def:physics:phyx-mini-0925:phyxminiproblems-problemphyxmini0925-tiltedsolenoidloopfigure}
  Independent apparatus data. In particular, loopMagneticFlux and all four areas are observables; none is defined from an answer choice or from the requested numerical flux
\end{definition}
\begin{proof}
  This is the declared model definition of Tilted Solenoid Loop Setup.
\end{proof}
\begin{definition}[Matches Problem Description]
  \label{def:physics:phyx-mini-0925:phyxminiproblems-problemphyxmini0925-matchesproblemdescription}
  \lean{PhyXMiniProblems.ProblemPhyXMini0925.MatchesProblemDescription}
  \uses{def:physics:phyx-mini-0925:phyxminiproblems-problemphyxmini0925-lengthincentimeters, def:physics:phyx-mini-0925:phyxminiproblems-problemphyxmini0925-magneticfluxdensityinteslas, def:physics:phyx-mini-0925:phyxminiproblems-problemphyxmini0925-tiltedsolenoidloopsetup}
  Numerical and qualitative information supplied by the problem prose
\end{definition}
\begin{proof}
  This is the declared model definition of Matches Problem Description.
\end{proof}
\begin{definition}[Matches Primary Tilt Figure]
  \label{def:physics:phyx-mini-0925:phyxminiproblems-problemphyxmini0925-matchesprimarytiltfigure}
  \lean{PhyXMiniProblems.ProblemPhyXMini0925.MatchesPrimaryTiltFigure}
  \uses{def:physics:phyx-mini-0925:phyxminiproblems-problemphyxmini0925-tiltedsolenoidloopsetup}
  Facts transcribed from the supplied raster rather than its noisy caption
\end{definition}
\begin{proof}
  This is the declared model definition of Matches Primary Tilt Figure.
\end{proof}
\begin{definition}[Has Physical Solenoid Loop Parameters]
  \label{def:physics:phyx-mini-0925:phyxminiproblems-problemphyxmini0925-hasphysicalsolenoidloopparameters}
  \lean{PhyXMiniProblems.ProblemPhyXMini0925.HasPhysicalSolenoidLoopParameters}
  \uses{def:physics:phyx-mini-0925:phyxminiproblems-problemphyxmini0925-lengthinmeters, def:physics:phyx-mini-0925:phyxminiproblems-problemphyxmini0925-magneticfluxdensityinteslas, def:physics:phyx-mini-0925:phyxminiproblems-problemphyxmini0925-tiltedsolenoidloopsetup}
  Positivity and fit conditions for the stated physical apparatus
\end{definition}
\begin{proof}
  This is the declared model definition of Has Physical Solenoid Loop Parameters.
\end{proof}
\begin{definition}[Satisfies Circular Area Geometry]
  \label{def:physics:phyx-mini-0925:phyxminiproblems-problemphyxmini0925-satisfiescircularareageometry}
  \lean{PhyXMiniProblems.ProblemPhyXMini0925.SatisfiesCircularAreaGeometry}
  \uses{def:physics:phyx-mini-0925:phyxminiproblems-problemphyxmini0925-lengthinmeters, def:physics:phyx-mini-0925:phyxminiproblems-problemphyxmini0925-areainsquaremeters, def:physics:phyx-mini-0925:phyxminiproblems-problemphyxmini0925-tiltedsolenoidloopsetup}
  Both stated diameters determine ordinary circular-disk areas
\end{definition}
\begin{proof}
  This is the declared model definition of Satisfies Circular Area Geometry.
\end{proof}
\begin{definition}[Satisfies Tilted Threading Geometry]
  \label{def:physics:phyx-mini-0925:phyxminiproblems-problemphyxmini0925-satisfiestiltedthreadinggeometry}
  \lean{PhyXMiniProblems.ProblemPhyXMini0925.SatisfiesTiltedThreadingGeometry}
  \uses{def:physics:phyx-mini-0925:phyxminiproblems-problemphyxmini0925-areainsquaremeters, def:physics:phyx-mini-0925:phyxminiproblems-problemphyxmini0925-degreestoradians, def:physics:phyx-mini-0925:phyxminiproblems-problemphyxmini0925-tiltedsolenoidloopsetup}
  The marked plane angle equals the angle between the solenoid axis and the loop area normal. Orthogonal projection multiplies the threaded surface area by the corresponding cosine. Centering and containment make that projection exactly the solenoid cross-section. These are geometric laws; they do not mention magnetic flux or an answer value
\end{definition}
\begin{proof}
  This is the declared model definition of Satisfies Tilted Threading Geometry.
\end{proof}
\begin{definition}[Models Uniform Confined Solenoid Field]
  \label{def:physics:phyx-mini-0925:phyxminiproblems-problemphyxmini0925-modelsuniformconfinedsolenoidfield}
  \lean{PhyXMiniProblems.ProblemPhyXMini0925.ModelsUniformConfinedSolenoidField}
  \uses{def:physics:phyx-mini-0925:phyxminiproblems-problemphyxmini0925-magneticfluxdensityinteslas, def:physics:phyx-mini-0925:phyxminiproblems-problemphyxmini0925-tiltedsolenoidloopsetup}
  Physlib's vector field is interpreted in tesla coordinates. It is uniform and axial throughout the solenoid interior and vanishes outside, faithfully encoding the field-confinement idealization responsible for using the small solenoid cross-section rather than the full loop area
\end{definition}
\begin{proof}
  This is the declared model definition of Models Uniform Confined Solenoid Field.
\end{proof}
\begin{definition}[Satisfies Uniform Magnetic Surface Flux Law]
  \label{def:physics:phyx-mini-0925:phyxminiproblems-problemphyxmini0925-satisfiesuniformmagneticsurfacefluxlaw}
  \lean{PhyXMiniProblems.ProblemPhyXMini0925.SatisfiesUniformMagneticSurfaceFluxLaw}
  \uses{def:physics:phyx-mini-0925:phyxminiproblems-problemphyxmini0925-areainsquaremeters, def:physics:phyx-mini-0925:phyxminiproblems-problemphyxmini0925-magneticfluxdensityinteslas, def:physics:phyx-mini-0925:phyxminiproblems-problemphyxmini0925-magneticfluxinwebers, def:physics:phyx-mini-0925:phyxminiproblems-problemphyxmini0925-tiltedsolenoidloopsetup}
  For a uniform field, the oriented surface flux is the field magnitude times the normal cosine times the actual threaded surface area. The law is stated before substituting either the projection geometry or any numerical data
\end{definition}
\begin{proof}
  This is the declared model definition of Satisfies Uniform Magnetic Surface Flux Law.
\end{proof}
\begin{definition}[Answer Choice]
  \label{def:physics:phyx-mini-0925:phyxminiproblems-problemphyxmini0925-answerchoice}
  \lean{PhyXMiniProblems.ProblemPhyXMini0925.AnswerChoice}
  Multiple-choice labels in the source problem
\end{definition}
\begin{proof}
  This is the declared model definition of Answer Choice.
\end{proof}
\begin{definition}[displayed Flux In Webers]
  \label{def:physics:phyx-mini-0925:phyxminiproblems-problemphyxmini0925-displayedfluxinwebers}
  \lean{PhyXMiniProblems.ProblemPhyXMini0925.displayedFluxInWebers}
  \uses{def:physics:phyx-mini-0925:phyxminiproblems-problemphyxmini0925-answerchoice}
  Magnetic-flux readout in webers printed beside each answer choice
\end{definition}
\begin{proof}
  This is the declared model definition of displayed Flux In Webers.
\end{proof}
\begin{definition}[displayed Flux Resolution In Webers]
  \label{def:physics:phyx-mini-0925:phyxminiproblems-problemphyxmini0925-displayedfluxresolutioninwebers}
  \lean{PhyXMiniProblems.ProblemPhyXMini0925.displayedFluxResolutionInWebers}
  \uses{def:physics:phyx-mini-0925:phyxminiproblems-problemphyxmini0925-answerchoice}
  Resolution of the final displayed digit of each two-significant-figure value
\end{definition}
\begin{proof}
  This is the declared model definition of displayed Flux Resolution In Webers.
\end{proof}
\begin{definition}[Rounds To Displayed Flux]
  \label{def:physics:phyx-mini-0925:phyxminiproblems-problemphyxmini0925-roundstodisplayedflux}
  \lean{PhyXMiniProblems.ProblemPhyXMini0925.RoundsToDisplayedFlux}
  \uses{def:physics:phyx-mini-0925:phyxminiproblems-problemphyxmini0925-magneticfluxinwebers, def:physics:phyx-mini-0925:phyxminiproblems-problemphyxmini0925-tiltedsolenoidloopsetup, def:physics:phyx-mini-0925:phyxminiproblems-problemphyxmini0925-answerchoice, def:physics:phyx-mini-0925:phyxminiproblems-problemphyxmini0925-displayedfluxinwebers, def:physics:phyx-mini-0925:phyxminiproblems-problemphyxmini0925-displayedfluxresolutioninwebers}
  The physical flux rounds to a displayed choice at its stated resolution
\end{definition}
\begin{proof}
  This is the declared model definition of Rounds To Displayed Flux.
\end{proof}
\begin{definition}[Is Closest Displayed Choice]
  \label{def:physics:phyx-mini-0925:phyxminiproblems-problemphyxmini0925-isclosestdisplayedchoice}
  \lean{PhyXMiniProblems.ProblemPhyXMini0925.IsClosestDisplayedChoice}
  \uses{def:physics:phyx-mini-0925:phyxminiproblems-problemphyxmini0925-magneticfluxinwebers, def:physics:phyx-mini-0925:phyxminiproblems-problemphyxmini0925-tiltedsolenoidloopsetup, def:physics:phyx-mini-0925:phyxminiproblems-problemphyxmini0925-answerchoice, def:physics:phyx-mini-0925:phyxminiproblems-problemphyxmini0925-displayedfluxinwebers}
  A selected answer is no farther from the physical flux than any other
\end{definition}
\begin{proof}
  This is the declared model definition of Is Closest Displayed Choice.
\end{proof}
\begin{definition}[recorded Dataset Answer]
  \label{def:physics:phyx-mini-0925:phyxminiproblems-problemphyxmini0925-recordeddatasetanswer}
  \lean{PhyXMiniProblems.ProblemPhyXMini0925.recordedDatasetAnswer}
  \uses{def:physics:phyx-mini-0925:phyxminiproblems-problemphyxmini0925-answerchoice}
  The answer label recorded in the dataset metadata
\end{definition}
\begin{proof}
  This is the declared model definition of recorded Dataset Answer.
\end{proof}
% --- Archon declaration topology end ---
% --- Archon physics formalization source end ---
